% SPDX-FileCopyrightText: © 2026 Idiap Research Institute <contact@idiap.ch>
%
% SPDX-FileContributor: Andre Freitas <andre.freitas@idiap.ch>
% Neuro-symbolic AI Group
%
% SPDX-FileContributor: Danilo Gusicuma <danilo.gusicuma@idiap.ch>
% Neuro-symbolic AI Group
%
% SPDX-License-Identifier: MIT
% Generated from step-explanations.json
\paragraph{Step-by-step reconstruction.}
Reviewers examine a correlation-decay argument whose calculation concerns a visible sector but whose conclusion concerns the full physical spectrum.
The following is a source-guided illustrative reconstruction. Internal procedures are prescriptions; the native declarations do not establish their execution.
\begin{enumerate}[leftmargin=*,label=\arabic*.]
\item \textbf{Ground the reviewed state} (\texttt{review\_\allowbreak{}state}).
\begin{description}[style=standard,labelindent=0pt,leftmargin=5.5em,labelwidth=5em,labelsep=0.5em,align=left,font=\normalfont\bfseries,itemsep=0.15em,parsep=0pt,topsep=0.3em]
\item[Purpose] Ground criticism and repair in the same state.
\item[Input] Evidence bundle; implemented argument; persisted claims
\item[Method] Separate mathematical results, implementation behavior, assumptions and open obligations before interpreting progress.
\item[Check] Disagreement about scope must remain explicit rather than averaged into confidence.
\item[Output] Formalized review state: visible decay; unsupported global extension
\item[Next] Both parallel branches consume this same state; their coexistence gives no independence guarantee.
\item[Example] Write the state as a valid-looking visible-sector calculation plus the additional, unsupported claim that every physical excitation is gapped.
\end{description}
\item \textbf{Localize the defeater} (\texttt{defeater\_\allowbreak{}assessment}).
\begin{description}[style=standard,labelindent=0pt,leftmargin=5.5em,labelwidth=5em,labelsep=0.5em,align=left,font=\normalfont\bfseries,itemsep=0.15em,parsep=0pt,topsep=0.3em]
\item[Purpose] Locate the unsupported implication.
\item[Input] Grounded review state, shared with repair formation
\item[Method] Challenge quantifiers and representation coverage, separating a defect in the inference from a defect in implementation.
\item[Check] A useful diagnosis names the missing burden and the claim it affects.
\item[Output] Defeater assessment: missing observable-density premise
\item[Next] The join later compares this finding with repairs; the repair branch does not consume it directly.
\item[Example] An abstract direct-sum spectral countermodel preserves the reported correlations while adding energies approaching zero. It defeats the unrestricted correlation-to-gap inference and identifies the missing completeness premise; it is not a constructed Yang--Mills theory.
\end{description}
\item \textbf{Propose recovery options} (\texttt{repair\_\allowbreak{}options}).
\begin{description}[style=standard,labelindent=0pt,leftmargin=5.5em,labelwidth=5em,labelsep=0.5em,align=left,font=\normalfont\bfseries,itemsep=0.15em,parsep=0pt,topsep=0.3em]
\item[Purpose] Propose repairs that preserve useful content.
\item[Input] Formalized review state; no parallel refutation output
\item[Method] Propose a changed representation, missing bridge or narrower claim, explaining what each option would preserve.
\item[Check] A repair proposal must alter the defective dependency rather than relabel the current result.
\item[Output] Repair alternatives: richer observables or sectorial theorem
\item[Next] Final calibration must test these options against the separately produced defeater.
\item[Example] From the visible/global mismatch, propose either constructing a dense observable-generated subspace or explicitly restricting the conclusion to the visible sector.
\end{description}
\item \textbf{Judge repairs against criticism} (\texttt{reorientation\_\allowbreak{}decision}).
\begin{description}[style=standard,labelindent=0pt,leftmargin=5.5em,labelwidth=5em,labelsep=0.5em,align=left,font=\normalfont\bfseries,itemsep=0.15em,parsep=0pt,topsep=0.3em]
\item[Purpose] Match repairs to diagnosed defects.
\item[Input] Defeater assessment + parallel repair options
\item[Method] Check which proposal addresses the missing implication and record whether understanding sharpened, stalled or regressed.
\item[Check] A proposed repair does not resolve a defect merely by appearing in the same packet.
\item[Output] Reorientation assessment: repair assignment or conceptual reorientation
\item[Next] Original policy routes assignments; native return does not execute repair or prove density.
\item[Example] The narrower claim resolves the overstatement; the richer observable construction remains a research task. More curves without new sectors do not answer the defeater.
\end{description}
\end{enumerate}
\paragraph{Changed-input contrast.}
If a proposed repair only adds more decay curves from the same incomplete observable family, the join should reject it as irrelevant to density. A narrower sectorial claim changes the conclusion honestly.
